\section{Proposed Methodology}
\label{sec:methodology}

\subsection{Approach overview}
\label{subsec:method-overview}

Each hybrid-cloud challenge identified in Chapter~\ref{Chapter1} is mapped to a concrete
module that addresses it. Table~\ref{tab:problem-module} makes this mapping explicit and
provides the traceability from problem to solution that motivates the remainder of the
section.

\begin{table}[ht]
\centering
\caption{Mapping from hybrid-cloud challenges to solution modules.}
\label{tab:problem-module}
\begin{tabular}{|L{5.2cm}|L{6.8cm}|}
\hline
\textbf{Challenge} & \textbf{Module / mechanism} \\
\hline
Weak security linkage between public and private cloud & One gate and one risk engine
applied to both the CDK and the Ansible side. \\
\hline
Absence of a clear, repeatable process & The synth--gate--decide--deploy pipeline with
fixed thresholds and audit records. \\
\hline
Fragmented tooling with inconsistent output & Scanner adapters that normalise findings
into a common schema, followed by cross-source deduplication. \\
\hline
Misconfiguration risk amplified by fast AI generation & The generate--gate--regenerate
feedback loop that injects findings back into the prompt. \\
\hline
Context-free prioritisation & A risk score that combines severity with cost, live
compliance, and a machine-learning code-risk signal. \\
\hline
\end{tabular}
\end{table}

\subsection{AI-assisted generation and the regeneration loop}
\label{subsec:method-generation}

Zone~1 generates synthesizable CDK Python code from a natural-language prompt. Code
generation is placed behind a provider abstraction so that two backends --- a managed
foundation-model service and a hosted inference API --- are interchangeable behind a
single provider flag and share an identical generation contract. Model defaults are
centralised and can be overridden explicitly.

Generation alone does not guarantee secure output, so it is wrapped in a closed loop.
The loop generates code, synthesizes it, gates it, and --- if the gate rejects ---
rebuilds the prompt with the machine-readable findings injected before retrying, up to a
bounded number of attempts. Each attempt is persisted for reproducibility (the prompt snapshot, the generated
code, the synthesis output, and the gate report), and the successful attempt is promoted
to a \texttt{passed/} folder.

Figure~\ref{fig:method-regen} presents the loop as a flowchart, emphasising the two
exit conditions: an accepting decision (\emph{pass} or \emph{review}) leaves the loop
with a gated application, whereas exhausting the attempt budget $N$ leaves it with a
failure that surfaces to the operator.

\begin{figure}[ht]
\centering
\begin{tikzpicture}[node distance=6mm and 12mm]
\node[proc] (gen) {Generate\\code};
\node[proc, right=of gen] (syn) {\texttt{cdk synth}};
\node[dec, right=of syn] (ok) {synth\\ok?};
\node[proc, below=30mm of ok] (gate) {Gate\\(score)};
\node[dec, left=of gate] (rej) {reject?};
\node[proc, left=of rej] (done) {Return\\gated app};
\node[dec, below=10mm of gen] (budget) {$i<N$?};
\draw[flow] (gen) -- (syn);
\draw[flow] (syn) -- (ok);
\draw[flow] (ok) -- node[right,font=\scriptsize]{yes} (gate);
\draw[flow] (gate) -- (rej);
\draw[flow] (rej) -- node[above,font=\scriptsize]{no} (done);
\draw[flow] (rej) |- (0,-3.5) -| (budget.south) node[pos=0.75,right,font=\scriptsize]{yes: inject findings};
\draw[flow] (ok.south) |- (budget.east) node[pos=0.6,above,font=\scriptsize]{no: synth error};
\draw[flow] (budget) -- node[right,font=\scriptsize]{yes} (gen);
\node[proc, left=8mm of budget] (fail) {Fail};
\draw[flow] (budget) -- node[above,font=\scriptsize]{no} (fail);
\end{tikzpicture}
\caption{Regeneration loop as a flowchart. The loop exits on an accepting decision or
when the attempt budget $N$ is exhausted.}
\label{fig:method-regen}
\end{figure}

\subsection{Pre-validation}
\label{subsec:method-prevalidation}

Before any security analysis, each side is validated syntactically. On the CDK side,
\texttt{cdk synth} must succeed (otherwise the loop treats the failure as feedback), and
\texttt{cdk diff} previews the change. On the Ansible side,
\texttt{ansible-playbook -{}-syntax-check} validates the plays and
\texttt{ansible-playbook -{}-check -{}-diff} performs a dry run. This staging
(\emph{validate $\rightarrow$ dry run $\rightarrow$ deploy}) is symmetric across the two
sides and ensures the gate only ever scores artifacts that are at least well formed.

\subsection{Security evaluation and cross-source deduplication}
\label{subsec:method-evaluation}

Zone~2 analyses the synthesized artifact with several independent scanners, each wrapped
in an adapter that normalises its output into a common finding schema with the fields
\texttt{severity}, \texttt{source}, \texttt{message}, \texttt{resource\_id},
\texttt{template}, and \texttt{category}. Because several tools frequently report the
same underlying issue (for example, an open administrative port may be flagged by both a
heuristic check and a policy scanner), the findings are deduplicated by the key
$(\texttt{resource\_id}, \texttt{template}, \texttt{category})$. When duplicates collide,
the highest severity is retained and the source labels are merged. Deduplication is
essential: without it, redundant reports would inflate the score and distort the
decision. Figure~\ref{fig:method-gate} shows the analysis pipeline: the synthesized
artifact fans out to the parallel scanner adapters, whose normalized findings are merged,
deduplicated, and passed to the risk-scoring engine.

\begin{figure}[ht]
\centering
\begin{tikzpicture}[node distance=5mm and 10mm]
\node[io] (tmpl) {Synthesized\\artifact};
% scanners stacked
\node[proc, right=10mm of tmpl, yshift=15mm] (s1) {Static / policy\\scanners};
\node[proc, below=of s1] (s2) {Cost \& live\\compliance};
\node[proc, below=of s2] (s3) {ML code-risk\\model};
\node[proc, right=5mm of s2] (norm) {Normalize\\to schema};
\node[proc, right=5mm of norm] (dedup) {Deduplicate\\(\emph{res,tmpl,cat})};
\node[dec, right=5mm of dedup] (engine) {Risk\\engine};
\draw[flow] (tmpl.east) -- ++(6mm,0) |- (s1.west);
\draw[flow] (tmpl.east) -- ++(6mm,0) |- (s2.west);
\draw[flow] (tmpl.east) -- ++(6mm,0) |- (s3.west);
\draw[flow] (s1.east) -| (norm.north);
\draw[flow] (s2) -- (norm);
\draw[flow] (s3.east) -| (norm.south);
\draw[flow] (norm) -- (dedup);
\draw[flow] (dedup) -- (engine);
\end{tikzpicture}
\caption{Security-evaluation pipeline: parallel scanners, normalization to a common
schema, cross-source deduplication, and aggregation by the risk-scoring engine.}
\label{fig:method-gate}
\end{figure}

\subsection{Risk-scoring engine (overview)}
\label{subsec:method-scoring}

The deduplicated findings, together with the cost, live-compliance, and
machine-learning code-risk signals, are aggregated by the gate into a single numeric
score that is then mapped to one of three decisions --- \emph{pass}, \emph{review}, or
\emph{reject}. Because this aggregation is the analytical heart of the proposed process,
it is presented in full in Section~\ref{sec:risk-scoring}, which develops the multi-factor
formulation, the additive scoring implementation, the machine-learning code-risk model,
the deduplication-driven explainability, and a worked example. Here it suffices to note
that the score is an \emph{additive} combination of weighted findings, banded cost, live
compliance violations, and a bounded machine-learning contribution, and that the decision
thresholds are fixed so that identical inputs always yield identical decisions on both
sides of the hybrid environment.

\subsection{Deployment governance}
\label{subsec:method-governance}

The decision governs deployment. A \emph{pass} deploys automatically. A \emph{review}
requires an explicit approval, which is recorded together with the approver's cloud
identity, the run identifier, the score, and the decision. A \emph{reject} never
deploys; when the code was generated, the regeneration loop is invoked, and otherwise a
rejection record is written for the author to act on. Approval and rejection records are
append-only, which yields the immutable audit trail required by the design goals.

\subsection{Monitoring and the operational loop}
\label{subsec:method-monitoring}

After a successful deployment, Zone~3 attaches post-deployment monitoring and closes an
operational loop. Each gate decision is persisted to SSM Parameter Store, published as an
EventBridge event, and emitted as CloudWatch metrics and structured logs. An
event-driven function reacts to compliance drift and rollback events, enabling automated
remediation. Because the on-premises nodes are registered as managed instances, their
compliance events flow through the same loop, so the entire hybrid environment is
observed uniformly rather than as two disconnected halves.
